% ==========================================================================
%  Chapter 74 -- Two-Way FE² Coupling: Homogenised Tangent and
%                Staggered Iteration
% ==========================================================================

\chapter{Two-Way FE\textsuperscript{2} Coupling}
\label{ch:two-way-fe2}

\paragraph{Normative status.}
This chapter is retained as a historical design record.  The current normative
behaviour of the public multiscale API is documented in
Chapter~\ref{ch:multiscale-publication-audit}.  Where terminology or behaviour
differs, Chapter~\ref{ch:multiscale-publication-audit} prevails.


% =========================================================================
\section{Introduction}
\label{sec:fe2-intro}

The one-way downscaling framework (Chapters~\ref{ch:multiscale_seismic}
--\ref{ch:nl-sub-model-evolver}) transfers kinematics from the global
beam model to local continuum sub-models, but does not feed the sub-model
response back into the beam analysis.  This chapter develops the
mathematical and computational framework for \emph{two-way}
FE\textsuperscript{2} coupling, where the homogenised sub-model tangent
modifies the beam section stiffness and the two scales are iterated until
equilibrium is reached at both levels simultaneously.


% =========================================================================
\section{FE\textsuperscript{2} Framework}
\label{sec:fe2-framework}

% -------------------------------------------------------------------------
\subsection{Scale Separation}

\textbf{Macro-scale.}
The global structural model consists of fibre-section beam elements
with $n_{\text{gp}}$ Gauss points along the axis.  At each section
Gauss point the generalised forces and deformations are:
\begin{equation}
  \label{eq:fe2-gen}
  \mathbf{s} = \{N, \, M_y, \, M_z, \, V_y, \, V_z, \, T\}^T, \qquad
  \mathbf{e} = \{\varepsilon_0, \, \kappa_y, \, \kappa_z,
                 \, \gamma_y, \, \gamma_z, \, \theta'\}^T.
\end{equation}

\textbf{Micro-scale.}
Each critical section Gauss point is associated with a prismatic RVE
(Representative Volume Element), modelled as a 3D continuum domain
$\Omega_{\text{RVE}}$ with optional embedded rebar
(Chapter~\ref{ch:mixed-model-rebar}).  The RVE is driven by
displacement boundary conditions derived from the macro-scale section
kinematics via the Timoshenko kinematic field
(Chapter~\ref{ch:structural-scale-reconstruction}).

% -------------------------------------------------------------------------
\subsection{Homogenisation: Stress to Section Forces}

The section forces are obtained by volume-averaging the RVE stress
field and integrating over the cross-section:
\begin{equation}
  \label{eq:fe2-homog-N}
  N = \frac{1}{L_{\text{RVE}}} \int_{\Omega_{\text{RVE}}}
      \sigma_{zz} \; \mathrm{d}\Omega,
\end{equation}
\begin{equation}
  \label{eq:fe2-homog-M}
  M_y = \frac{1}{L_{\text{RVE}}} \int_{\Omega_{\text{RVE}}}
        \sigma_{zz} \, x \; \mathrm{d}\Omega, \quad
  M_z = -\frac{1}{L_{\text{RVE}}} \int_{\Omega_{\text{RVE}}}
        \sigma_{zz} \, y \; \mathrm{d}\Omega,
\end{equation}
\begin{equation}
  \label{eq:fe2-homog-V}
  V_y = \frac{1}{L_{\text{RVE}}} \int_{\Omega_{\text{RVE}}}
        \sigma_{yz} \; \mathrm{d}\Omega, \quad
  V_z = \frac{1}{L_{\text{RVE}}} \int_{\Omega_{\text{RVE}}}
        \sigma_{xz} \; \mathrm{d}\Omega,
\end{equation}
where $z$ is the beam axis direction and $L_{\text{RVE}}$ is the RVE
length.  These integrals reduce to discrete summations over Gauss
points with appropriate weights.

% -------------------------------------------------------------------------
\subsection{Section Forces from Boundary Reactions}
\label{sec:fe2-reactions}

An alternative---and more robust---approach extracts the section
resultants directly from the assembled internal-force vector.
After convergence of the RVE finite-element solve with displacement
$\mathbf{u}$, the global internal-force vector is
\begin{equation}
  \mathbf{f}_{\mathrm{int}} = \bigcup_{e}
    \int_{\Omega_e} \mathbf{B}^T \boldsymbol{\sigma}(\mathbf{u})
    \; \mathrm{d}\Omega.
\end{equation}
At the constrained Face~B boundary nodes, $\mathbf{f}_{\mathrm{int}}$
gives the reaction forces.  The section resultants follow from
summation over all Face~B nodes~$i$:
\begin{equation}
  \mathbf{F} = \sum_{i \in \text{Face\,B}} \mathbf{f}_i,
  \qquad
  \mathbf{M} = \sum_{i \in \text{Face\,B}}
    (\mathbf{x}_i - \mathbf{c}) \times \mathbf{f}_i,
\end{equation}
where $\mathbf{c}$ is the section centroid.  Rotating to the beam
local frame via $\mathbf{R}$ gives $N$, $V_y$, $V_z$ from
$\mathbf{R}\mathbf{F}$ and $M_t$, $M_y$, $M_z$ from
$\mathbf{R}\mathbf{M}$.

This approach is preferred for the perturbation tangent loop because
it evaluates the constitutive response from the displacement field
at each perturbation step, avoiding reliance on stored material state
which may be stale between non-committed SNES solves.


% =========================================================================
\section{Homogenised Section Tangent}
\label{sec:fe2-tangent}

% -------------------------------------------------------------------------
\subsection{Perturbation-Based Tangent Computation}

The $6 \times 6$ section tangent matrix
$\mathbf{D}_{\text{hom}}$ relates infinitesimal increments of
section deformations to section forces:
\begin{equation}
  \label{eq:fe2-Dhom}
  \delta\mathbf{s} = \mathbf{D}_{\text{hom}} \, \delta\mathbf{e},
  \qquad
  D_{ij} = \frac{\partial s_i}{\partial e_j}.
\end{equation}

Since the sub-model solver provides only the nonlinear response
$\mathbf{s}(\mathbf{e})$, the tangent is computed by numerical
(central) differentiation.  For each section deformation component
$e_j$, $j = 1, \ldots, 6$:
\begin{enumerate}
  \item Perturb: $\mathbf{e}^+ = \mathbf{e} + h \, \mathbf{e}_j$,
    $\mathbf{e}^- = \mathbf{e} - h \, \mathbf{e}_j$.
  \item Solve two nonlinear sub-models (or apply the kinematic field
    and re-homogenise without full Newton solve for small $h$).
  \item Differentiate:
    $\mathbf{D}_{\text{hom}}(:, j) \approx
    (\mathbf{s}^+ - \mathbf{s}^-) / (2h)$.
\end{enumerate}

The perturbation amplitude $h$ should be chosen small enough for
accuracy but large enough to avoid round-off:
\begin{equation}
  h = \max(\epsilon_r \, |e_j|, \; h_{\min}),
\end{equation}
where $\epsilon_r \approx 10^{-6}$ and $h_{\min} \approx 10^{-10}$.

For efficiency, the forward-difference approximation may be used when
$\mathbf{s}(\mathbf{e})$ is already available at the current state:
\begin{equation}
  \mathbf{D}_{\text{hom}}(:, j) \approx
  \frac{\mathbf{s}(\mathbf{e} + h\,\mathbf{e}_j) - \mathbf{s}(\mathbf{e})}{h}.
\end{equation}
This halves the number of sub-model evaluations at the cost of
first-order accuracy.

% -------------------------------------------------------------------------
\subsection{Structure of $\mathbf{D}_{\text{hom}}$}

For a rectangular RC section under combined axial--flexural loading, the
tangent has the block structure:
\begin{equation}
  \mathbf{D}_{\text{hom}} =
  \begin{bmatrix}
    EA     & ES_y   & ES_z   & 0      & 0      & 0 \\
    ES_y   & EI_y   & EI_{yz}& 0      & 0      & 0 \\
    ES_z   & EI_{yz}& EI_z   & 0      & 0      & 0 \\
    0      & 0      & 0      & GA_y   & 0      & 0 \\
    0      & 0      & 0      & 0      & GA_z   & 0 \\
    0      & 0      & 0      & 0      & 0      & GJ
  \end{bmatrix},
\end{equation}
where the axial--flexural block ($3 \times 3$) is fully populated
(coupling between $N$, $M_y$, $M_z$) and the shear--torsion block
($3 \times 3$) is diagonal for symmetric sections.  After cracking, the
off-diagonal axial--flexural coupling terms grow as the neutral axis
shifts.


% =========================================================================
\section{Stiffness Injection into the Beam Element}
\label{sec:fe2-injection}

The homogenised tangent $\mathbf{D}_{\text{hom}}$ replaces the
fibre-based section tangent at the critical Gauss point.  Two
injection strategies are considered:

\paragraph{Strategy~A: Section tangent override.}
The beam element's section stiffness at the critical Gauss point is
replaced wholesale:
\begin{equation}
  \mathbf{C}_t^{(k)} \leftarrow \mathbf{D}_{\text{hom}}^{(k)},
  \qquad k \in \mathcal{G}_{\text{crit}},
\end{equation}
where $\mathcal{G}_{\text{crit}}$ is the set of critical Gauss points.
The element stiffness is then re-integrated:
\begin{equation}
  \mathbf{K}_e = \int_0^L
    \mathbf{B}^T \, \mathbf{C}_t(s) \, \mathbf{B} \; \mathrm{d}s,
\end{equation}
where $\mathbf{C}_t(s)$ uses the fibre section tangent at non-critical
Gauss points and $\mathbf{D}_{\text{hom}}$ at critical ones.

\paragraph{Strategy~B: Homogenised material section adapter.}
A \texttt{HomogenizedMaterialSection} class that wraps the
$6 \times 6$ tangent as a ``virtual fibre section'' and returns
$\mathbf{D}_{\text{hom}}$ from \texttt{compute\_section\_tangent()}.
This approach integrates cleanly with the existing beam element
template, which expects sections that satisfy the
\texttt{BeamSectionConcept}.


% =========================================================================
\section{Staggered Iteration Scheme}
\label{sec:fe2-stagger}

The FE\textsuperscript{2} coupling is driven by a staggered iteration
loop within each dynamic time step:

\begin{enumerate}
  \item \textbf{Global solve}: Solve the macro-scale dynamic analysis
    for one time step, producing updated beam kinematics.
  \item \textbf{Local solve}: For each critical element, update the
    RVE boundary conditions and solve the nonlinear sub-model.
  \item \textbf{Homogenise}: Extract effective section forces and
    tangent from each sub-model via volume averaging and numerical
    differentiation.
  \item \textbf{Inject}: Replace the beam section tangent at critical
    Gauss points with $\mathbf{D}_{\text{hom}}$.
  \item \textbf{Convergence check}: Compare the sub-model section
    forces with the beam section forces at the critical Gauss points:
    \begin{equation}
      \label{eq:fe2-conv}
      \frac{\|\mathbf{s}_{\text{beam}} - \mathbf{s}_{\text{sub}}\|}
           {\|\mathbf{s}_{\text{sub}}\| + \epsilon_0}
      < \epsilon_{\text{stag}}.
    \end{equation}
  \item If not converged, return to step~1 with the updated section
    stiffness.
\end{enumerate}

\paragraph{Convergence properties.}
The staggered scheme converges when the spectral radius of the
coupling operator is less than unity.  For quasi-static problems
with moderate nonlinearity, 2--4 staggered iterations typically
suffice.  Under dynamic loading with explicit time integration at
the macro-scale, a single staggered pass per step may be adequate
if the time step is sufficiently small.


% =========================================================================
\section{Planned Implementation Architecture}
\label{sec:fe2-architecture}

The two-way coupling builds on three new components and one
modification to the existing codebase:

\begin{enumerate}
  \item \textbf{\texttt{homogenize\_tangent()}} in
    \texttt{HomogenizedSection.hh}: Performs 6 (or 12) perturbation
    solves on the sub-model to compute $\mathbf{D}_{\text{hom}}$
    via central (or forward) differences.  Returns a
    $6 \times 6$ Eigen matrix.

  \item \textbf{\texttt{HomogenizedMaterialSection}}: An adapter class
    satisfying \texttt{BeamSectionConcept} that stores the $6 \times 6$
    tangent and returns it from the section stiffness query.  This
    allows seamless integration with \texttt{BeamElement}'s Gauss-point
    integration loop.

  \item \textbf{Stiffness injection in \texttt{BeamElement}}: A method
    \texttt{override\_section\_tangent(gp\_index, D\_hom)} that replaces
    the fibre-based tangent at a specific Gauss point with the
    homogenised tangent.  The beam element re-integrates using the
    overridden section for subsequent stiffness matrix and internal
    force computations.

  \item \textbf{Staggered iteration loop}: Implemented in the main
    driver (e.g.\ \texttt{main\_lshaped\_multiscale.cpp}) as an outer
    loop wrapping the dynamic solver step and the sub-model updates.
    The loop tracks convergence via \eqref{eq:fe2-conv} and exits
    after a maximum number of staggered passes.
\end{enumerate}


% =========================================================================
\section{Summary}
\label{sec:fe2-summary}

The two-way FE\textsuperscript{2} coupling closes the multiscale loop:
\begin{itemize}
  \item \textbf{Downward}: Beam kinematics $\to$ RVE boundary
    conditions $\to$ nonlinear sub-model solve.
  \item \textbf{Upward}: Sub-model stress $\to$ homogenised section
    forces and tangent $\to$ beam section stiffness override.
  \item \textbf{Iteration}: Staggered scheme iterates macro/micro
    solves until equilibrium at both scales.
\end{itemize}

Combined with the arc-length solver (Chapter~\ref{ch:arc-length}) and
mixedmodel reinforcement (Chapter~\ref{ch:mixed-model-rebar}), this
framework enables full nonlinear multiscale analysis of
reinforced-concrete structures through limit and softening regimes.
